\section{Spritzgießen und Faserorientierung}
\label{sec:spritzgiessen}

Zahnräder aus thermoplastischen Kunststoffen werden vorwiegend durch Spritzgießen hergestellt; daneben kommen das Gießen und spanende Verfahren zum Einsatz \cite[S.~25]{VDI2736-1}. Welches Verfahren gewählt wird, hängt vom Zahnradwerkstoff, von der geforderten Verzahnungsqualität sowie von der Stückzahl und von Form und Größe der Räder ab. Bei hohen Stückzahlen ist das Spritzgießen in der Regel sehr kostengünstig und bietet eine große Gestaltungsfreiheit. Werden besonders hohe Anforderungen an die Verzahnungsqualität gestellt, kann ein spanendes Nachbearbeiten erforderlich sein \cite[S.~25]{VDI2736-1}. Spritzgegossene Kunststoffzahnräder weisen gegenüber Stahlrädern eine um den Faktor zehn bis hundert kleinere Steifigkeit auf, abhängig von der Faserverstärkung, dazu schlechtere thermische Eigenschaften und eine geringere Genauigkeit \cite[S.~1]{FVA856-I}. Dem stehen das geringere Gewicht, die Trockenlauffähigkeit und die Eignung zur kostengünstigen Massenproduktion gegenüber \cite[S.~1]{FVA856-I}.

Beim Spritzgießen wird ein thermoplastisches Granulat aufgeschmolzen und unter hohem Druck in die Kavität eines Werkzeugs gespritzt, in der die Schmelze beim Erstarren die Form der Kavität annimmt \cite[S.~25]{VDI2736-1}. Bei teilkristallinen Kunststoffen hängt die Verarbeitungsschwindung unter anderem von der Wanddicke ab. Teile mit großen Wanddicken erfordern längere Zykluszeiten und neigen zu Lunkern und Verzug, weshalb in der Regel nur Zahnräder mit Wanddicken bis $10$\,mm spritzgegossen werden \cite[S.~25]{VDI2736-1}. Glasfaserverstärkte Kunststoffe schwinden weniger, die Glasfaserorientierung ruft jedoch von der Fließrichtung abhängige Schwindungsunterschiede und damit Verzug hervor \cite[S.~25]{VDI2736-1}. Als Richtwerte nennt die Richtlinie für glasfaserverstärkte Polyamide eine Verarbeitungsschwindung von $0{,}2$\,\% bei stark behinderter und $0{,}8$\,\% bei freier Schwindung, für unverstärkte Polyamide $1{,}2$\,\% bis $2{,}2$\,\% \cite[S.~26]{VDI2736-1}. Die Schwindung wird hauptsächlich von der Massetemperatur, der Temperatur des Formnests, dem Massedruck, dem Anspritzpunkt, der Einspritzgeschwindigkeit und der Nachdruckzeit beeinflusst \cite[S.~26]{VDI2736-1}. Schwankungen dieser Verarbeitungsgrößen verursachen Maßänderungen am Zahnrad. Die Nachschwindung hängt von Höhe und Gleichmäßigkeit der Temperatur in der Kavität ab und lässt sich durch höhere Werkzeugtemperaturen so weit vermindern, dass ein nachfolgendes Tempern entfallen kann \cite[S.~26 f.]{VDI2736-1}.

Für eine hohe Fertigungsqualität verlangt die Richtlinie vom Werkzeug möglichst gleiche thermische Bedingungen im gesamten Bauteil und ein balanciertes Angusssystem bei Mehrfachwerkzeugen \cite[S.~27 f.]{VDI2736-1}. Hinzu kommen ein radial gleichmäßiger Druckverlauf im Formnest und ein verzugsarmes Auswerfen. Die negative Verzahnung im Werkzeugeinsatz wird durch Fräsen, Stoßen, Kalt-Einsenken, Erodieren oder galvanoplastisches Abformen hergestellt. Schwindungsabhängige Maß- und Formabweichungen des Zahnrads sind bereits beim Fertigen der Einsätze zu berücksichtigen. Ist die Schwindung der Zahnform nicht genau bekannt, sind mehrere Korrekturschritte am Durchmesser wie an Zahnform und Zahnhöhe notwendig \cite[S.~28]{VDI2736-1}. Weil die Geometrie der Verzahnung nur einmal bei der Erstellung des Werkzeugs gefertigt wird, lassen sich durch Spritzgießen auch komplexe Zahngeometrien herstellen. Hinterschnitte bedingen Schieber im Werkzeug, die dessen Genauigkeit und damit die der Bauteile verringern \cite[S.~28]{VDI2736-1}. Diese Gestaltungsfreiheit gegenüber der spanenden Fertigung ist auch der Ansatz, mit dem Kassem et al. die Fußrundung von Kunststoffzahnrädern frei gestalten \cite[S.~2]{Kas22}.

Unterschiedliche Wanddicken bewirken Schwindungsunterschiede, die innere Spannungen hervorrufen und zu Verzug führen. Aussparungen im Radkörper sind deshalb möglichst symmetrisch anzuordnen, und ein tulpenförmiger Querschnitt mit gegeneinander versetztem Teller und Nabe ist nach Möglichkeit zu vermeiden \cite[S.~37 f.]{VDI2736-1}. Die Lage der Anbindung ist ein Kompromiss zwischen Rundheitsgenauigkeit und Herstellkosten. Für eine möglichst hohe Rundheitsgenauigkeit soll das Formteil zentral angebunden werden; in vielen Fällen hat sich der 3-Punkt-Kaltkanalanschnitt bewährt, der im 3-Platten-Werkzeug beim Öffnen abgetrennt wird \cite[S.~38]{VDI2736-1}. Bei faserverstärkten Werkstoffen ist im Bereich der Bindenähte mit Rundheitsabweichungen zu rechnen \cite[S.~38]{VDI2736-1}. Alle Bauteilbereiche sollen bis zu ihrem Einfrieren mit Schmelze versorgt werden, wozu die Wanddicken zum Anschnitt hin etwas zunehmen müssen. Da kein Auslauf für ein Verzahnungswerkzeug benötigt wird, lassen sich die Zähne an eine Planfläche anbinden und dadurch verstärken \cite[S.~38]{VDI2736-1}. Rippen verbessern die Steifigkeit des Radkörpers, bergen aber durch Werkstoffanhäufungen an den Anschlussstellen die Gefahr zusätzlichen Verzugs sowie von Rundheitsabweichungen und Einfallstellen \cite[S.~39]{VDI2736-1}. Nach Illenberger hat sich in der Praxis das Schirmangussverfahren bewährt, bei kleineren Rädern auch ein Mehrfach-Punktanguss auf der Stirnseite; Versteifungsrippen sollen speichenförmig von der Mitte nach außen geführt werden \cite[S.~5 f.]{Ill22}. Im Mehrkomponentenspritzguss können Radkörper und Welle-Nabe-Verbindung aus einem steifen faserverstärkten Werkstoff mit geringer Schwindung, die Verzahnung dagegen aus einem tribologisch optimierten Werkstoff gespritzt werden \cite[S.~6]{Ill22}. VDI 2736 Blatt 1 nennt dafür das Zwei-Komponenten-Verfahren mit unterschiedlichen Werkstoffen an Zahnflanke und Zahnfuß \cite[S.~14]{VDI2736-1}.

Die Fertigungstoleranzen sind erst nach sorgfältiger Prüfung der technischen Anforderungen festzulegen, weil zu enge Toleranzen das Verfahren unwirtschaftlich machen \cite[S.~30]{VDI2736-1}. Im Betrieb können Abweichungen aus Zahnverformung sowie aus Temperatur- und Feuchteänderung des Kunststoffs so groß werden, dass sie gegenüber den Fertigungstoleranzen dominieren. Kopfkreisdurchmesser sowie Zahndicke und Zahnweite werden bei spritzgegossenen Rädern durch die Spritzbedingungen beeinflusst. Die Lageabweichung des Zahnkranzes gegenüber der Drehachse hängt von der Genauigkeit des Werkzeugs, von Schwindungsunterschieden innerhalb des Radkörpers und vom Verzug ab \cite[S.~30 f.]{VDI2736-1}. Weil sich Kunststoffzahnräder leicht verformen und sich kleinere Verzahnungsabweichungen beim Abwälzen ausgleichen, werden größere Abweichungen zugelassen als bei Rädern aus Metall. Die Qualität wird ausschließlich über das Ergebnis der Wälzprüfung festgelegt, für Flankenlinien- und Zahndickenabweichungen können bis zu zwei Qualitäten mehr zugelassen werden \cite[S.~33 f.]{VDI2736-1}. Als Erfahrungswerte ordnet die Richtlinie Getrieben mit hohen Anforderungen die Qualitäten 8 bis 9 und Getrieben mit normalen Anforderungen die Qualitäten 10 bis 11 nach DIN 3963 zu \cite[S.~34]{VDI2736-1}. Gemessen wird frühestens $24$\,h nach der Herstellung im Normalklima; Grundlage der Bewertung ist der arithmetische Mittelwert aus Messungen an mehr als drei Zähnen je Flanke \cite[S.~33]{VDI2736-1}.

Bei kurzfaserverstärkten Thermoplasten richtet die Strömung während der Füllphase die Fasern aus, und diese Orientierung ist nicht zufällig \cite[S.~280]{Sto18}. Fasern in einem Strömungsbereich mit deutlichem Geschwindigkeitsgradienten erfahren so lange eine Drehung, bis sich beide Faserenden mit gleicher Geschwindigkeit bewegen. Dieser Effekt ist im wandnahen Bereich durch die höheren Schergeschwindigkeiten deutlich stärker ausgeprägt als in der Fließkanalmitte, in der sich die Fasern translatorisch mit der Schmelze bewegen \cite[S.~281]{Sto18}. Dehnströmungen quer zur Fließrichtung an der Schmelzefront und bei einer sich verbreiternden Fließfront richten die Fasern dagegen quer aus. Die Randschicht weist deshalb immer eine deutliche Vorzugsrichtung in Fließrichtung auf, während der Orientierungsgrad der Mittelschicht vom Füllverhalten des Formteils abhängt \cite[S.~281]{Sto18}. Wie dick die Randschicht im Verhältnis zur Mittelschicht ist, folgt aus dem Geschwindigkeitsprofil der Strömung. Nimmt die Viskosität mit der Schergeschwindigkeit stark ab, bildet sich eine Blockströmung mit dünner, hochorientierter Randschicht und breiter Mittelschicht \cite[S.~281]{Sto18}. Erhard beschreibt für eine Rechteckplatte mit Bandanguss einen Schichtenaufbau aus faserarmem Randfilm, dünner regelloser Schicht und einer Scherschicht mit Fasern überwiegend in Fließrichtung \cite[S.~255 f.]{Erh08}. Die Fasern der Kernschicht orientieren sich wegen der dort herrschenden Zugspannungen überwiegend quer zur Fließrichtung. Aus den strömungsbedingten Faserorientierungen resultieren erhebliche Festigkeits- und Steifigkeitsanisotropien, und an Bindenähten stellt sich eine strenge Längsorientierung ein \cite[S.~257 und 259]{Erh08}.

Die lokale Faserorientierung hängt damit direkt vom Füllbild und von den Eigenschaften der Schmelze ab \cite[S.~282]{Sto18}. Für ein Formteil ist sie nur über eine Prozesssimulation zu ermitteln, sofern keine CT-Aufnahme gefertigter Teile vorliegt. Einzelne Fasern beschreibt die Spritzgießsimulation nicht. Stattdessen schätzen mathematische Modelle die lokale Orientierungsverteilung aus dem Geschwindigkeitsfeld der Strömung ab; bekannt sind das Folgar-Tucker-Modell und dessen Erweiterung, das RSC-Modell \cite[S.~282]{Sto18}. Beide Ansätze setzen eine homogene Faserverteilung und gleichbleibende Faserlängen voraus und erfassen keine Rückwirkung der Fasern auf die Strömung. In Wirklichkeit werden Fasern entlang des Fließwegs fortlaufend gebrochen \cite[S.~283]{Sto18}. Ergebnis der Rechnung ist in jedem Element des Füllmodells der Orientierungstensor zweiter Stufe $a_{\mathrm{ij}}$, aus dem sich Eigenwerte und Eigenvektoren bilden lassen \cite[S.~282 f.]{Sto18}. Die Füllsimulation der Prüfräder dieser Arbeit wurde mit CADMOULD gerechnet (\secref{sec:fe_modell}). \Note{PH: CADMOULD-Bilder der Füllung und Orientierung.}

Im Vorhaben FVA 856 I wurde die Faserorientierung an Zugproben und Zahnrädern aus PA46-GF30 mit Computertomographie bei einer Auflösung von $4$\,\textmu m bestimmt \cite[S.~15]{FVA856-I}. Am Zahnrad wurde ein Stück aus drei Zähnen gescannt. Die Auswertung liefert je Zone zweidimensionale Projektionen der Faserorientierung in den drei Ebenen als Histogramm und Farbdarstellung sowie den Orientierungstensor mit seinen Eigenwerten und Eigenvektoren und den Anisotropie-Index \cite[S.~15]{FVA856-I}. Eine Projektion in eine Ebene macht die Vorzugsorientierung schnell deutlich, ihre Ausprägung lässt sich aus der Breite des Peaks aber nur schwer abschätzen. Im Orientierungstensor liegt die vollständige räumliche Information; er bezieht sich auf das globale Koordinatensystem des CT-Scans, und seine Eigenwerte und Eigenvektoren beschreiben die Faserorientierung im Hauptsystem der Zone \cite[S.~17]{FVA856-I}. Aus den Eigenwerten folgt der Anisotropie-Index
\begin{equation}
    \mathrm{AI} = 1 - \frac{\lambda_{\mathrm{min}}}{\lambda_{\mathrm{max}}}
    \label{eq:anisotropie_index}
\end{equation}
mit dem kleinsten Eigenwert $\lambda_{\mathrm{min}}$ und dem größten Eigenwert $\lambda_{\mathrm{max}}$ \cite[S.~17]{FVA856-I}. Der Index liegt zwischen $0$ für die ideal isotrope Orientierung, bei der alle Eigenwerte $1/3$ betragen, und $1$ für die ideal anisotrope Orientierung mit einem einzigen Eigenwert von $1$. Zonen mit verschiedener Vorzugsrichtung lassen sich damit in der Ausprägung ihrer Ausrichtung vergleichen \cite[S.~17]{FVA856-I}. \figref{fig:quelle_fva856_bild16} zeigt diese Auswertung an einer Zone am linken Zahnfuß.

\begin{figure}[htbp]
    \centering
    \includegraphics[width=0.8\textwidth]{kap2/quelle_fva856_bild16}
    \caption{Auswertung der Faserorientierung in einer Zone am linken Zahnfuß aus der Computertomographie: Projektion in eine Ebene, Orientierungstensor mit Eigenwerten und Eigenvektoren, nach \cite[S.~18]{FVA856-I}.}
    \label{fig:quelle_fva856_bild16}
\end{figure}

Die Prüfverzahnung des Vorhabens hat den Achsabstand $60$\,mm und den Normalmodul $2$\,mm bei je 30 Zähnen \cite[S.~104 f.]{FVA856-I}. Das Ritzel besteht aus einsatzgehärtetem und geschliffenem 16MnCr5 mit der Zahnbreite $15$\,mm, das Rad aus spritzgegossenem PA46-GF30 mit der Zahnbreite $12$\,mm; beide wurden von ZF Friedrichshafen bereitgestellt \cite[S.~104 f.]{FVA856-I}. Die Räder bestehen aus demselben Werkstoff wie die Zugproben, einem mit 30 Gewichtsprozent Kurzglasfasern verstärkten Polyamid 46 der Bezeichnung TW300F6 \cite[S.~10 und 105]{FVA856-I}. Schliffbilder zeigten an allen Zähnen eine vergleichbare Faserorientierung, sodass die einzeln im CT untersuchten Zähne als repräsentativ gelten \cite[S.~15]{FVA856-I}. Im Zahninneren auf Höhe des Zahnfußes verblieben Poren aus dem Spritzgießprozess, die wegen ihrer innenliegenden Lage nicht als versagensauslösend eingeschätzt wurden. Die Vorzugsorientierung folgt im CT erwartungsgemäß der Zahnoberfläche, bedingt durch die Fließrichtung beim Spritzgießen \cite[S.~16]{FVA856-I}. Der Zahn wurde in drei Arten von Zonen eingeteilt \cite[S.~18]{FVA856-I}. Oberflächenzonen folgen der Faserflussrichtung, je Seite das Tal, der Fuß sowie die untere und die obere Flanke, dazu die Spitze. Im Zahninneren liegen die Zonen Mitte und unterer Bereich, und Übergangszonen vermitteln dazwischen. \figref{fig:quelle_fva856_bild17} fasst die Zonen mit ihren Eigenwerten und dem Anisotropie-Index zusammen. Die am stärksten ausgeprägte Anisotropie liegt in den Zahnflanken und im Zahnfußbereich vor, das Innere des Zahns ist deutlich isotroper orientiert, und die Übergangszonen liegen dazwischen \cite[S.~18]{FVA856-I}. Zwischen rechten und linken Zonen stimmen die Eigenwerte nahezu überein \cite[S.~19]{FVA856-I}. Für die FE-Simulation wurde daraus ein Zonenplan mit der Vorzugsorientierung in der Frontebene abgeleitet, in dem die drei Übergangszonen zu einer zusammengefasst sind (\figref{fig:quelle_fva856_bild18}) \cite[S.~19]{FVA856-I}. Im FE-Modell wurde die Verzahnung in 34 Segmente mit je einem lokalen Koordinatensystem aufgeteilt, dessen 1-Richtung der Faservorzugsorientierung entspricht. Die Segmentgröße ist auf den Normalmodul bezogen, damit andere Baugrößen gleich segmentiert werden können \cite[S.~53 f.]{FVA856-I}.

\begin{figure}[htbp]
    \centering
    \includegraphics[width=0.85\textwidth]{kap2/quelle_fva856_bild17}
    \caption{Zonen des Zahns mit den Eigenwerten des Orientierungstensors und dem Anisotropie-Index aus der Computertomographie; die Eigenwerte der rechten Zonen sind nahezu gleich denen der linken, nach \cite[S.~19]{FVA856-I}.}
    \label{fig:quelle_fva856_bild17}
\end{figure}

\begin{figure}[htbp]
    \centering
    \includegraphics[width=0.5\textwidth]{kap2/quelle_fva856_bild18}
    \caption{Der FE-Analyse übergebene Zonenkarte mit der Faservorzugsrichtung in der Frontebene, nach \cite[S.~19]{FVA856-I}.}
    \label{fig:quelle_fva856_bild18}
\end{figure}

Mikromechanische Rechnungen mit den gemessenen Faserlängen und den Eigenwerten der Zonen ergaben einen Modul in Faservorzugsrichtung von $4{,}38$\,GPa in der Zone Mitte bis $6{,}41$\,GPa in der Zone linker Fuß \cite[S.~23]{FVA856-I}. Quer dazu reicht der Modul von $3{,}21$\,GPa bis $3{,}70$\,GPa. Stärker anisotrope Zonen erzeugen erwartungsgemäß einen höheren Modul in Vorzugsrichtung und einen geringeren quer dazu; die außen liegenden Zonen weisen damit den höchsten Modul auf. Über die Zahnbreite gedrittelt liegt der Modul im mittleren Drittel der äußeren Zonen um etwa $5$\,\% bis $15$\,\% niedriger, in den Übergangszonen um bis zu $20$\,\% \cite[S.~24]{FVA856-I}. Als Grund wird eine stärkere Verwirbelung der Fasern in diesem Drittel vermutet. Die Fasern im Zahnrad sind länger als in den Zugproben, deren Verteilung durch Faserbrüche einen wesentlich höheren Anteil sehr kurzer Fasern unter $100$\,\textmu m enthält \cite[S.~20]{FVA856-I}. Der Faserdurchmesser betrug in allen Proben etwa $10$\,\textmu m. Deshalb hält der Bericht fest, dass die im Zugversuch ermittelten Kennwerte nicht ohne Weiteres auf die Zahnradsimulation übertragbar sind \cite[S.~24 f.]{FVA856-I}.

Die Kette von der Prozesssimulation über die Struktur zur Beanspruchung ist das Prinzip, nach dem kurzfaserverstärkte Bauteile berechnet werden. Programme zur Prozess-Struktur-Kopplung wie CONVERSE übertragen die lokale Faserinformation automatisch in die Struktursimulation; über den Orientierungstensor werden die Hauptorientierungsrichtung und die Orientierungsverteilung an jeder Stelle des Formteils beschrieben \cite[S.~86]{Sto18}. Stommel et al. nennen drei Wege, diese Information zu nutzen \cite[S.~87 f.]{Sto18}. Der einfachste übernimmt allein die Hauptrichtung als lokale 1-Richtung mit transversal-isotropen Kennwerten. Daneben stehen eine zweistufige Homogenisierung über Faserklassen und die Mittelung der Steifigkeit nach Advani und Tucker mit den Orientierungstensoren zweiter und vierter Stufe. Der erste Weg ist ausreichend genau, wenn die Orientierungsverteilung hinreichend eng ist, und kann ohne Kopplungssoftware durch manuell angenommene Vorzugsorientierungen im kritischen Bauteilbereich ersetzt werden \cite[S.~87]{Sto18}. Kassem et al. bilden die Faserorientierung ohne Prozesssimulation nach. Aus den Tangentenvektoren der Zahngeometrie wird für den Schwerpunkt jedes Elements eine Faserrichtung interpoliert, und jedes Element erhält ein lokales Koordinatensystem mit der 1-Richtung in Faserrichtung \cite[S.~10]{Kas22}. Die Übereinstimmung mit dem CT wird als sehr gut bewertet. Für die elastischen Konstanten der hochorientierten Randschicht greifen sie auf die mikromechanische Rechnung des Vorhabens zurück, weil der Dreischichtaufbau der Zugproben die Übertragbarkeit der gemessenen Kennwerte verfälscht \cite[S.~11]{Kas22}. Das Modell teilt das Rad in kleine Elemente und weist jedem eine vom Faserwinkel abhängige Werkstoffbeschreibung zu \cite[S.~2593]{Kas23}. Im Multiphysik-Modell des Vorhabens ergeben sich höhere Zahnfußspannungen als mit einem isotropen linear elastischen Modell, was neben der Dehnratenabhängigkeit auf die Inhomogenität des Werkstoffs zurückgeführt wird \cite[S.~128]{FVA856-I}. In Schichten mit höherem Modul treten höhere Spannungen auf, und bei einem Biegebalken nehmen die äußeren Bereiche die Last hauptsächlich auf. Die größte Zahnfußspannung wird dort in der Randschicht der Zugseite ausgewertet, weil Zahnfußbrüche an dieser Stelle beginnen \cite[S.~129]{FVA856-I}. Auf dieser Grundlage leiten Kassem et al. mit der klassischen Laminattheorie eine analytische Biegespannung für den Zahnfuß her \cite{Kas24}. Der aus dem CT abgeleitete Faserwinkel geht darin als stetige Funktion über die Zahndicke ein. Das Verfahren ist so gefasst, dass es den Formfaktor der VDI 2736 ersetzen kann, und stimmt mit der FE-Rechnung deutlich besser überein als die isotrope Rechnung \cite{Kas24}. Der Orientierungsgrad, der nahe der Oberfläche deutlich höher ist als im Zahninneren und die Zahnfußspannung weiter erhöht, bleibt darin wie in der FE-Validierung unberücksichtigt \cite{Kas24}. Auch der Abschlussbericht sieht eine analytische Ableitung der Zahnfußspannung faserverstärkter Räder erst nach Berücksichtigung des Orientierungsgrads als möglich an \cite[S.~122]{FVA856-I}. Wie weit die Orientierung aus der eigenen Füllsimulation den CT-Zonen des Vorhabens entspricht, wird in Kapitel 5 gezeigt.

% Herkunft:
% Absatz 1 (Herstellverfahren): VDI 2736 Blatt 1 S. 25 (Abschnitt 7 und 7.1: vorwiegend Spritzgießen, Gießen
%   und Spanen, Auswahlkriterien, kostengünstig bei hohen Stückzahlen, Gestaltungsfreiheit, spanendes
%   Nachbearbeiten); FVA 856 I S. 1 (Kapitel 2: Steifigkeit Faktor zehn bis hundert kleiner, thermische
%   Eigenschaften, Genauigkeit; Vorteile Gewicht, Trockenlauf, Dämpfung, Korrosion, Massenproduktion).
% Absatz 2 (Verfahren, Schwindung): VDI 2736 Blatt 1 S. 25 (7.1 und 7.1.1: Granulat, Kavität, Wanddicke,
%   Lunker und Verzug, 10 mm, GF schwindet weniger, fließrichtungsabhängige Schwindung), S. 26 (Einflussgrößen
%   der Schwindung, Maßänderungen, Nachschwindung; Tabelle 6: PA 66-GF X 0,2/0,8/0,8 %, PA 6 1,2/1,8/2,0 %,
%   PA 66 1,2/1,8/2,2 %), S. 26 f. (Tempern entfällt).
% Absatz 3 (Werkzeug): VDI 2736 Blatt 1 S. 27 f. (7.1.2 Anforderungen), S. 28 (Herstellung der negativen
%   Verzahnung, Korrekturschritte, komplexe Zahngeometrien, Hinterschnitte); Kassem 2022 S. 2 (design freedom
%   of injection molding, Abstract und Einleitung).
% Absatz 4 (Radkörper, Anguss): VDI 2736 Blatt 1 S. 37 f. (9.1.1: Wanddicken, Aussparungen, Tulpenform Bild
%   16), S. 38 (Anbindung, 3-Punkt-Kaltkanalanschnitt, Bindenähte bei faserverstärkten Werkstoffen, Schmelze
%   bis zum Einfrieren, Zähne an Planfläche), S. 39 (Rippen); Illenberger 2022 S. 5 f. (2.1.3 Schirmanguss,
%   Mehrfach-Punktanguss, speichenförmige Rippen, Mehrkomponentenspritzguss); VDI 2736 Blatt 1 S. 14
%   (2-K-Verfahren, Flanke z. B. POM/PTFE, Fuß z. B. POM-GF).
% Absatz 5 (Genauigkeit): VDI 2736 Blatt 1 S. 30 (Abschnitt 8, 8.1.1), S. 30 f. (8.1.3 Lageabweichung),
%   S. 33 f. (8.3.2 erreichbare Qualitäten, Wälzprüfung, zwei Qualitäten mehr, Tabelle 7), S. 33 (8.2: 24 h,
%   Normalklima, mehr als drei Zähne, Mittelwert).
% Absatz 6 (Rand- und Mittelschicht): Stommel 2018 S. 280 (5.2.1, Bild 5.8), S. 281 (Drehung der Fasern,
%   Dehnströmung, Randschicht, Blockströmung); Erhard 2008 S. 255 f. (7.1, Bild 7.6 Schichtenaufbau), S. 257
%   (Bindenaht mit Längsorientierung), S. 259 (Festigkeits- und Steifigkeitsanisotropien).
% Absatz 7 (Prozesssimulation): Stommel 2018 S. 282 (Füllbild, Folgar-Tucker Gl. (5.12), RSC Gl. (5.13)),
%   S. 283 (Annahmen, Faserschädigung, Eigenwerte lambda_p des Orientierungstensors), S. 282 f.; CADMOULD:
%   Entwurf 3.6 (sec:fe_modell, tab:cadmould).
% Absatz 8 (Tensor, AI): FVA 856 I S. 15 (6.2.1: nanotom 4 µm, drei Zähne, VG STUDIO MAX, Ergebnisse),
%   S. 17 (Projektion, Orientierungstensor, Eigenwerte, Gl. (1) AI = 1 - lambda_min/lambda_max, Grenzen 0 und 1).
% Formel (eq:anisotropie_index): FVA 856 I Gl. (1), S. 17.
% Absatz 9 (Prüfverzahnung, Zonen): FVA 856 I S. 104 f. (9.4, Tabelle 15: a = 60 mm, m_n = 2 mm, z = 30/30,
%   b = 15/12 mm, 16MnCr5, PA46-GF30, ZF), S. 10 (6.1: TW300F6, 30 Gew.-%, Konzelmann), S. 105 (gleicher
%   Werkstoff wie die Prüfkörper), S. 15 (Schliffbilder, Poren, Bild 13), S. 16 (6.2.1.2 Vorzugsorientierung
%   folgt der Oberfläche), S. 18 (drei Zonenarten, Bild 17 Verweis, Anisotropie je Zone), S. 19 (rechte Zonen,
%   Zonenplan Bild 18), S. 53 f. (8.1.3: 34 Segmente, lokales Koordinatensystem, Normalmodul).
% Absatz 10 (Steifigkeit der Zonen, Faserlänge): FVA 856 I S. 23 (6.3.2: E1 4,38 bis 6,41 GPa, E2 3,21 bis
%   3,70 GPa), S. 24 (Drittel: 5 bis 15 %, bis 20 %, Verwirbelung), S. 20 (6.2.2: Faserlänge, < 100 µm,
%   Durchmesser 10 ± 1 µm), S. 24 f. (6.3.3 Übertragbarkeit nicht automatisch gegeben).
% Absatz 11 (Prozess-Struktur-Beanspruchung): Stommel 2018 S. 86 (CONVERSE, Orientierungstensor), S. 87 f.
%   (drei Wege, Advani-Tucker Gl. (2.77)); Kassem 2022 S. 10 (Tangentenvektoren, Interpolation, lokales
%   Koordinatensystem, Figure 8), S. 11 (Dreischichtaufbau, Table 3); Kassem 2023 S. 2593 (Elemente mit
%   faserwinkelabhängiger Werkstoffbeschreibung); FVA 856 I S. 128 (11.2.3.1: Dehnrate, Inhomogenität,
%   Laminattheorie, Bild 143), S. 129 (Randschicht der Zugseite), S. 122 (11.2: keine analytische Ableitung,
%   Orientierungsgrad); Kassem/Müller/Oehler/Koch, Analytical tooth root stress calculation (Manuskript ohne
%   Seitenzahlen: Abstract, Kapitel 3 Gl. (5), Kapitel 5 Gl. (28) und (29), Kapitel 7 und 8); Verweis auf
%   Kapitel 5 (Abschnitt 5.1 Faserorientierung aus der Spritzgießsimulation, ADR-119).
% Bildwunsch:
% fig:quelle_fva856_bild16: FVA856-I.txt / KST_856_I_AB_1491_unlocked.pdf, Druckseite 18 (PDF-Seite 32),
%   Bild 16 "Exemplarische Ergebnisse der Faserorientierungsanalyse in einer Zone am linken Zahnfuß";
%   ganzes Bild ausschneiden (Projektion links, Tensor und Eigenwerte rechts). Datei kap2/quelle_fva856_bild16.
% fig:quelle_fva856_bild17: Druckseite 19 (PDF-Seite 33), Bild 17 "Definierte Zonen im Zahnrad und deren
%   Eigenwerte aus dem Orientierungstensor mit Anisotropie-Index"; ganzes Bild. Datei kap2/quelle_fva856_bild17.
% fig:quelle_fva856_bild18: Druckseite 19 (PDF-Seite 33), Bild 18 "Der FE-Analyse übergebene Zonenkarte mit
%   Faservorzugsrichtung in der Frontebene"; ganzes Bild, klein (0,5 Textbreite). Datei kap2/quelle_fva856_bild18.
%   Bild 17 und 18 könnten auch als ein zweiteiliges Bild (a) und (b) gesetzt werden, dann eine Umgebung weniger.
% Optional (nicht im Text eingebunden): Stommel 2018 Bild 5.8 "Entstehung von Faserorientierungen", S. 280
%   (PDF-Seite 298), als Prinzipbild der Rand- und Mittelschicht "nach [Sto18]".
% Platzhalter: \Note{PH: CADMOULD-Bilder der Füllung und Orientierung.} in Absatz 7.
% Bibwunsch: Erh08, Kas24 (Manuskript, Tagung und Jahr offen) in <SCRATCH>\bib_kap2_2_4.bib; FVA856-I, VDI2736-1,
%   Sto18, Kas22, Kas23, Ill22 stehen in litbib.bib.
% Nomenklatur:
% AI, -, Anisotropie-Index (aus den Eigenwerten des Orientierungstensors), FVA 856 I S. 17 und Formelzeichen
%   S. 5 ("AI Anisotropie-Index -"); nicht in Nomenklatur.tex.
% lambda_min, -, kleinster Eigenwert des Orientierungstensors, FVA 856 I S. 17; nicht in Nomenklatur.tex.
% lambda_max, -, größter Eigenwert des Orientierungstensors, FVA 856 I S. 17; nicht in Nomenklatur.tex
%   (Nomenklatur des Berichts S. 6: "lambda Eigenwert -").
% a_ij, -, Orientierungstensor zweiter Stufe, Stommel 2018 S. 282 ("aij: Orientierungstensor"), S. 88
%   ("Orientierungstensor 2. Stufe"); in Nomenklatur.tex als "Eigenwerte des Faserorientierungstensors
%   (Index 11 größter, 22 zweitgrößter Eigenwert)" vorhanden; Vorschlag: Benennung auf "Faserorientierungstensor
%   zweiter Stufe; im Hauptachsensystem stehen die Eigenwerte auf der Diagonalen (11 größter, 22 zweitgrößter)"
%   erweitern, damit Text und Werkstoffkarte dasselbe Zeichen führen.
% Offen:
% BA 2.1 Zeilen 777 bis 786 ([Sta20]: Kosten, Geräusch, Reibung; Hochleistungskunststoffe gezielt im Fuß und
%   an der Flanke): auf FVA 856 I S. 1 und VDI 2736 Blatt 1 S. 14 (2-K-Verfahren) umgestellt.
% BA Zeilen 794 bis 802 (3-Komponenten-Zahnrad IMS Gear, Bild 2.1, PA66 GF50/PEEK): gestrichen, keine Quelle;
%   die Aussage "Faserverstärkung für die Zahnflanke weniger geeignet, für den Zahnfuß positiv [Ill22]" steht
%   mit dem tatsächlichen Befund (Illenberger S. 106) in Abschnitt 2.4.
% BA Zeilen 803 bis 806 ([Bon14]: Verzug und Schwindung, radiale Versteifungsrippen): auf VDI 2736 Blatt 1
%   S. 25, 37 bis 39 und Illenberger S. 5 f. umgestellt.
% BA Zeilen 807 bis 812 ([İmr09]: unterschiedliche Zahnbreiten über die Profillinie): gestrichen, keine Quelle.
% Kas24: das Manuskript trägt weder Jahr noch Tagung; PDF-Metadaten "Manuskriptvorlage VDI Wissensforum",
%   Dateidatum 2024-12-04; Tagung, Jahr und VDI-Bericht vom Nutzer zu ergänzen, bis dahin Zitat ohne Seite.
% FVA 856 I Tabelle 15 nennt für den Werkstoff der Prüfräder nur "PA46-GF30" (ZF); die Bezeichnung TW300F6
%   steht auf S. 10 für das Granulat der Zugproben, S. 105 sagt "derselbe PA46-GF30 Werkstoff".
